\documentclass{article}

\usepackage[margin=1.2cm, top=1cm, bottom=1.8cm]{geometry}
\pagestyle{plain}
\linespread{1.4}

\usepackage{fontspec}
\usepackage{unicode-math}
\usepackage[main=greek, english]{babel}

\setmainfont{CMU Serif}
\setsansfont{CMU Sans Serif}
\setmonofont{CMU Typewriter Text}
\setmathfont{Latin Modern Math}

\usepackage{amsmath, enumitem}
\usepackage{dirtytalk, caption, subcaption, array, graphicx}
\usepackage{titlesec}
\usepackage{soul}
\usepackage[unicode]{hyperref}
\usepackage{verbatim}

\usepackage{pdfpages}

\usepackage{tikz}
\usetikzlibrary{arrows.meta, positioning}

\hypersetup{
    hidelinks,
    bookmarksopen=true,
    bookmarksnumbered=true,
    pdftitle={Συστήματα Αναμονής - 6η Εργαστηριακή Άσκηση},
    pdfauthor={Πέππας Μιχαήλ-Αθανάσιος}
}

\newcolumntype{P}[1]{>{\centering\arraybackslash}p{#1}}
\titleformat*{\subsubsection}{\large\bfseries}

\begin{document}

\begin{center}
    {
    \Large
    \textbf{ΕΘΝΙΚΟ ΜΕΤΣΟΒΙΟ ΠΟΛΥΤΕΧΝΕΙΟ\\
    ΣΧΟΛΗ ΗΛΕΚΤΡΟΛΟΓΩΝ ΜΗΧΑΝΙΚΩΝ\\
    ΚΑΙ ΜΗΧΑΝΙΚΩΝ ΥΠΟΛΟΓΙΣΤΩΝ}

    \vspace{1.5cm}
    \rule{18cm}{0.5pt}

    Τομέας Επικοινωνιών, Ηλεκτρονικής και Συστημάτων Πληροφορικής\\
    Εργαστήριο Διαχείρισης και Βέλτιστου Σχεδιασμού Δικτύων Τηλεματικής - NETMODE\\

    \rule{18cm}{0.5pt}

    \textbf{Συστήματα Αναμονής}\\
    \textbf{$\mathbf{6^{η}}$ Εργαστηριακή Άσκηση}\\
    }
\end{center}

\vspace{0.8cm}
\begin{center}
\includegraphics[width=7cm]{EMP Logo.png}
\end{center}

\Large
{
\vspace{0.8cm}
\noindent
\textbf{Στοιχεία Φοιτητή}
\begin{itemize}
    \item Όνομα: Πέππας Μιχαήλ-Αθανάσιος
    \item Αριθμός Μητρώου: 03121026
    \item $4^{\eta}$ Εργαστηριακή Ομάδα (Δευτέρα, 12:45-14:30)
    \item Ημερομηνία Διεξαγωγής Εργαστηρίου: 04.05.2026
    \item Ημερομηνία Παράδοσης Αναφοράς: 17.06.2026
\end{itemize}
}

\newpage
\tableofcontents

\newpage
\section{Πρόβλεψη χρονοσειράς φόρτου}

\subsection{Στόχος της άσκησης}

Στην παρούσα εργαστηριακή άσκηση μελετάμε το πρόβλημα της πρόβλεψης χρονοσειράς, με εφαρμογή σε δεδομένα ζήτησης υπολογιστικών πόρων GPU. Η λογική του προβλήματος συνδέεται άμεσα με τη σχεδίαση και αξιολόγηση συστημάτων αναμονής: πριν διαστασιολογηθεί ένα σύστημα, πρέπει να υπάρχει εκτίμηση όχι μόνο της φύσης του φόρτου, αλλά και της έντασής του στο μέλλον. Η πρόβλεψη φόρτου είναι επομένως απαραίτητη για capacity planning, επιλογή πόρων, αποφυγή συμφόρησης και αποφυγή υπερβολικής υπερδιαστασιολόγησης.

Το αρχείο δεδομένων περιέχει μία χρονοσειρά ωριαίων παρατηρήσεων. Σύμφωνα με την εκφώνηση, χρησιμοποιούμε τις πρώτες $1200$ ώρες ως ιστορικά δεδομένα εκπαίδευσης και προβλέπουμε τις αμέσως επόμενες ώρες. Στο συγκεκριμένο αρχείο \texttt{dataset.mat} που χρησιμοποιήθηκε στην εκτέλεση υπήρχαν συνολικά $1247$ παρατηρήσεις, άρα ο τελικός ορίζοντας πρόβλεψης ήταν οι ώρες $1201$ έως $1247$, δηλαδή $47$ μελλοντικές τιμές.

Το βασικό μέρος της άσκησης απαιτεί την εφαρμογή και σύγκριση δύο κατηγοριών μεθόδων:
\begin{itemize}
    \item μοντέλα γραμμικής παλινδρόμησης,
    \item μοντέλα πρόβλεψης με φίλτρα.
\end{itemize}

Το προαιρετικό μέρος της άσκησης, δηλαδή ο διαγωνισμός, δεν υλοποιήθηκε. Η παρούσα αναφορά αφορά μόνο το βασικό μέρος.

\subsection{Περιγραφή των δεδομένων}

Η χρονοσειρά συμβολίζεται με $Y_t$, όπου $t$ είναι ο διακριτός χρόνος σε ώρες. Τα δεδομένα χωρίστηκαν ως εξής:

\begin{center}
\begin{tabular}{|P{5cm}|P{5cm}|P{5cm}|}
\hline
\textbf{Τμήμα δεδομένων} & \textbf{Ώρες} & \textbf{Πλήθος παρατηρήσεων} \\
\hline
Training set & $1,\ldots,1200$ & $1200$ \\
\hline
Forecast/Test set & $1201,\ldots,1247$ & $47$ \\
\hline
Σύνολο & $1,\ldots,1247$ & $1247$ \\
\hline
\end{tabular}
\end{center}

Στο Σχήμα~\ref{fig:dataset_split} φαίνεται η συνολική χρονοσειρά και το σημείο διαχωρισμού μεταξύ ιστορικών και μελλοντικών τιμών.

\begin{figure}[htbp]
    \centering
    \includegraphics[width=0.95\textwidth]{lab6_1.png}
    \caption{Δεδομένα χρονοσειράς και διαχωρισμός training/test.}
    \label{fig:dataset_split}
\end{figure}

Παρατηρούμε ότι η χρονοσειρά παρουσιάζει έντονη μεταβλητότητα και ορισμένες μεγάλες αιχμές. Οι αιχμές αυτές είναι σημαντικές, επειδή δεν αποτελούν ομαλές περιοδικές συνιστώσες και επομένως δεν αναμένεται να προβλεφθούν με ακρίβεια από απλά ντιτερμινιστικά μοντέλα. Αντίθετα, τα μοντέλα μπορούν κυρίως να συλλάβουν τη μέση δομή, την τάση και την εποχικότητα της σειράς.

\section{Θεωρητικό υπόβαθρο}

\subsection{Η έννοια της πρόβλεψης}

Στη θεωρία πρόβλεψης χρονοσειρών, αν έχουν παρατηρηθεί οι τιμές:
\[
Y_1,Y_2,\ldots,Y_t,
\]
τότε η πλήρης πρόβλεψη για μια μελλοντική τιμή $Y_{t+\ell}$ είναι η υπό συνθήκη κατανομή:
\[
P(Y_{t+\ell} \mid Y_1,\ldots,Y_t).
\]
Η σημειακή πρόβλεψη είναι συνήθως η υπό συνθήκη μέση τιμή:
\[
\hat{Y}_{t+\ell} = E(Y_{t+\ell} \mid Y_1,\ldots,Y_t).
\]

Ωστόσο, στην πράξη δεν αρκεί μόνο μία σημειακή πρόβλεψη. Επειδή η μελλοντική τιμή είναι τυχαία και δεν μπορεί να προβλεφθεί χωρίς σφάλμα, πρέπει να συνοδεύεται από διάστημα πρόβλεψης. Για επίπεδο εμπιστοσύνης $95\%$, ένα prediction interval έχει γενικά τη μορφή:
\[
[\hat{Y}_{t+\ell} - c_{\ell},\hat{Y}_{t+\ell} + c_{\ell}],
\]
όπου το $c_{\ell}$ εκφράζει την αβεβαιότητα της πρόβλεψης στον ορίζοντα $\ell$.

Είναι σημαντικό να τονιστεί ότι τα prediction intervals δεν είναι ίδια με τα confidence intervals της μέσης τιμής. Τα prediction intervals αφορούν μελλοντικές πραγματικές παρατηρήσεις και επομένως είναι συνήθως ευρύτερα.

\subsection{Γραμμική παλινδρόμηση για πρόβλεψη}

Στη γραμμική παλινδρόμηση θεωρούμε ότι η χρονοσειρά γράφεται ως:
\[
Y_t = f_t(\boldsymbol{\beta}) + \epsilon_t,
\]
όπου $f_t(\boldsymbol{\beta})$ είναι γνωστή συνάρτηση ως προς τον χρόνο και γραμμική ως προς το διάνυσμα παραμέτρων $\boldsymbol{\beta}$, ενώ $\epsilon_t$ είναι ο τυχαίος όρος σφάλματος.

Στην παρούσα άσκηση χρησιμοποιήθηκαν όροι τάσης και περιοδικότητας. Για την περιοδικότητα χρησιμοποιήθηκαν ημιτονοειδείς και συνημιτονοειδείς όροι:
\[
\cos\left(\frac{2\pi t}{T}\right), \qquad
\sin\left(\frac{2\pi t}{T}\right),
\]
όπου $T$ είναι η περίοδος σε ώρες.

Επειδή τα δεδομένα είναι ωριαία, οι φυσικές περίοδοι που εξετάζονται είναι:
\[
T=24
\]
για ημερήσια περιοδικότητα, και:
\[
T=168 = 7 \cdot 24
\]
για εβδομαδιαία περιοδικότητα.

Για παράδειγμα, το μοντέλο που τελικά επιλέχθηκε στην κατηγορία της παλινδρόμησης έχει μορφή:
\[
Y_t = \beta_0 + \beta_1 t
+ \beta_2 \cos\left(\frac{2\pi t}{24}\right)
+ \beta_3 \sin\left(\frac{2\pi t}{24}\right)
+ \beta_4 \cos\left(\frac{2\pi t}{168}\right)
+ \beta_5 \sin\left(\frac{2\pi t}{168}\right)
+ \epsilon_t.
\]
Η εκτίμηση των παραμέτρων γίνεται με τη μέθοδο των ελαχίστων τετραγώνων. Αν θεωρήσουμε ότι τα σφάλματα είναι ανεξάρτητα, κανονικά και με σταθερή διασπορά, τότε η μέθοδος ελαχίστων τετραγώνων αντιστοιχεί και σε εκτίμηση μέγιστης πιθανοφάνειας.

\subsection{Φίλτρα και διαφοροποίηση χρονοσειράς}

Η δεύτερη οικογένεια μεθόδων βασίζεται στη χρήση φίλτρων. Ένα φίλτρο εφαρμόζεται στη χρονοσειρά ώστε να αφαιρέσει συστηματικές συνιστώσες, όπως τάση ή εποχικότητα, και να μετατρέψει το υπόλοιπο σε πιο απλή διαδικασία προς πρόβλεψη.

Χρησιμοποιούμε τον τελεστή οπισθοδρόμησης $B$, με:
\[
BY_t = Y_{t-1}.
\]
Ο τελεστής διαφοροποίησης τάξης $s$ είναι:
\[
\Delta_s = 1 - B^s.
\]
Άρα:
\[
\Delta_s Y_t = Y_t - Y_{t-s}.
\]

Για ημερήσια εποχικότητα χρησιμοποιείται:
\[
\Delta_{24}Y_t = Y_t - Y_{t-24},
\]
ενώ για εβδομαδιαία εποχικότητα:
\[
\Delta_{168}Y_t = Y_t - Y_{t-168}.
\]

Γενικά, αν το φίλτρο είναι:
\[
F(B)Y_t = c + \epsilon_t,
\]
τότε η πρόβλεψη γίνεται αντιστρέφοντας το φίλτρο. Στην υλοποίηση, το $c$ εκτιμήθηκε ως ο εμπειρικός μέσος της φιλτραρισμένης σειράς στο training set. Αν ο μέσος αυτός είναι κοντά στο μηδέν, τότε το μοντέλο ανάγεται στην κλασική υπόθεση ότι η φιλτραρισμένη σειρά είναι θόρυβος μηδενικής μέσης τιμής.

\subsection{Μετρικές αξιολόγησης}

Για τη σύγκριση των μοντέλων χρησιμοποιήθηκαν οι εξής μετρικές:

\begin{itemize}
    \item Mean Absolute Error:
    \[
    MAE = \frac{1}{m}\sum_{i=1}^{m}|Y_i-\hat{Y}_i|.
    \]
    \item Root Mean Squared Error:
    \[
    RMSE = \sqrt{\frac{1}{m}\sum_{i=1}^{m}(Y_i-\hat{Y}_i)^2}.
    \]
    \item Coverage των $95\%$ prediction intervals:
    \[
    Coverage = \frac{\#\{Y_i \in [L_i,U_i]\}}{m}.
    \]
    \item Median CRPS.
\end{itemize}

Το CRPS χρησιμοποιήθηκε ως βασικό κριτήριο επιλογής του καλύτερου μοντέλου ανά κατηγορία, επειδή αξιολογεί ταυτόχρονα τη θέση της προβλεπτικής κατανομής και το εύρος της αβεβαιότητας. Για κανονική προβλεπτική κατανομή με μέση τιμή $\mu$ και τυπική απόκλιση $\sigma$, το CRPS υπολογίζεται από:
\[
CRPS = \sigma\left[
a(2\Phi(a)-1)+2\phi(a)-\frac{1}{\sqrt{\pi}}
\right],
\]
όπου:
\[
a = \frac{y-\mu}{\sigma}.
\]
Εδώ $\Phi$ και $\phi$ είναι αντίστοιχα η αθροιστική συνάρτηση κατανομής και η πυκνότητα της τυπικής κανονικής κατανομής.

\section{Μεθοδολογία μοντελοποίησης}

\subsection{Επιλογή υποψήφιων μοντέλων}

Για να μην βασιστούμε σε ένα αυθαίρετα επιλεγμένο μοντέλο, δοκιμάστηκαν πέντε παραλλαγές για κάθε κατηγορία μεθόδων. Δεν χρησιμοποιήθηκε ξεχωριστό validation set, επειδή η εκφώνηση ζητά να χρησιμοποιηθούν οι πρώτες $1200$ ώρες ως ιστορικό και να προβλεφθούν οι επόμενες ώρες. Έτσι, όλα τα μοντέλα προσαρμόστηκαν στο ίδιο training set και αξιολογήθηκαν στο ζητούμενο forecast/test interval.

Τα πέντε υποψήφια μοντέλα γραμμικής παλινδρόμησης ήταν:

\begin{center}
\begin{tabular}{|P{2.4cm}|P{10.5cm}|P{3.2cm}|}
\hline
\textbf{Μοντέλο} & \textbf{Περιγραφή} & \textbf{Λογική} \\
\hline
LR1 & Σταθερός όρος + ημερήσια περιοδικότητα $T=24$ & Βασικό ημερήσιο μοντέλο \\
\hline
LR2 & Γραμμική τάση + ημερήσια περιοδικότητα $T=24$ & Προσθήκη αργής μεταβολής \\
\hline
LR3 & Γραμμική τάση + ημερήσια $T=24$ + εβδομαδιαία $T=168$ περιοδικότητα & Ημερήσια και εβδομαδιαία δομή \\
\hline
LR4 & Τετραγωνική τάση + ημερήσια $T=24$ + εβδομαδιαία $T=168$ περιοδικότητα & Έλεγχος καμπυλότητας \\
\hline
LR5 & Γραμμική τάση + περίοδοι $24$, $12$, $168$ & Πρόσθετη ημιημερήσια αρμονική \\
\hline
\end{tabular}
\end{center}

Τα πέντε υποψήφια μοντέλα φίλτρων ήταν:

\begin{center}
\begin{tabular}{|P{2.4cm}|P{5cm}|P{8.5cm}|}
\hline
\textbf{Μοντέλο} & \textbf{Φίλτρο} & \textbf{Λογική} \\
\hline
F1 & $\Delta_{24}$ & Αφαίρεση ημερήσιας εποχικότητας \\
\hline
F2 & $\Delta_1\Delta_{24}$ & Αφαίρεση τάσης και ημερήσιας εποχικότητας \\
\hline
F3 & $\Delta_{168}$ & Αφαίρεση εβδομαδιαίας εποχικότητας \\
\hline
F4 & $\Delta_1\Delta_{168}$ & Αφαίρεση τάσης και εβδομαδιαίας εποχικότητας \\
\hline
F5 & $\Delta_{24}\Delta_{168}$ & Συνδυασμός ημερήσιας και εβδομαδιαίας διαφοροποίησης \\
\hline
\end{tabular}
\end{center}

Η επιλογή του καλύτερου μοντέλου σε κάθε κατηγορία έγινε με βάση το μικρότερο median CRPS. Η επιλογή αυτή είναι κατάλληλη, επειδή δεν ενδιαφέρει μόνο το μέσο τετραγωνικό σφάλμα της σημειακής πρόβλεψης, αλλά και το αν τα prediction intervals είναι χρήσιμα και όχι υπερβολικά στενά ή υπερβολικά πλατιά.

\section{Αποτελέσματα γραμμικής παλινδρόμησης}

\subsection{Σύγκριση υποψήφιων μοντέλων}

Τα αποτελέσματα των πέντε μοντέλων παλινδρόμησης φαίνονται στον Πίνακα~\ref{tab:lr_results}.

\begin{table}[htbp]
\centering
\small
\caption{Αποτελέσματα υποψήφιων μοντέλων γραμμικής παλινδρόμησης.}
\label{tab:lr_results}
\begin{tabular}{|P{7cm}|P{1.4cm}|P{2.1cm}|P{1.7cm}|P{1.7cm}|P{2cm}|}
\hline
\textbf{Μοντέλο} & \textbf{Παράμ.} & \textbf{Train MSE} & \textbf{MAE} & \textbf{RMSE} & \textbf{Median CRPS} \\
\hline
LR1: constant + daily 24 & 3 & 521883.168 & 538.035 & 791.283 & 260.300 \\
\hline
LR2: linear trend + daily 24 & 4 & 513540.892 & 458.657 & 710.656 & 232.094 \\
\hline
LR3: linear trend + daily 24 + weekly 168 & 6 & 498776.282 & 419.878 & 675.712 & 209.130 \\
\hline
LR4: quadratic trend + daily 24 + weekly 168 & 7 & 497531.338 & 437.483 & 702.799 & 215.583 \\
\hline
LR5: linear trend + daily 24 + half-daily 12 + weekly 168 & 8 & 490046.575 & 402.893 & 644.796 & 235.488 \\
\hline
\end{tabular}
\end{table}

Το καλύτερο μοντέλο με βάση το median CRPS είναι το:
\[
LR3:\quad \text{linear trend + daily 24 + weekly 168}.
\]
Παρατηρούμε ότι το $LR5$ έχει μικρότερο $MAE$ και $RMSE$ από το $LR3$, όμως έχει χειρότερο median CRPS. Αυτό σημαίνει ότι, παρότι οι σημειακές προβλέψεις του $LR5$ είναι ελαφρώς καλύτερες ως προς τα απόλυτα και τετραγωνικά σφάλματα, η συνολική πιθανοτική πρόβλεψη, δηλαδή ο συνδυασμός σημειακής πρόβλεψης και αβεβαιότητας, είναι καλύτερη για το $LR3$.

Η παρατήρηση αυτή είναι σημαντική σε σχέση με το overfitting. Το $LR5$ έχει περισσότερες παραμέτρους και μικρότερο training MSE, αλλά αυτό δεν μεταφράζεται σε καλύτερο probabilistic forecast. Αντίθετα, το $LR3$ αποτελεί πιο ισορροπημένη επιλογή, επειδή περιέχει τις δύο βασικές εποχικότητες χωρίς να προσθέτει υπερβολική πολυπλοκότητα.

\subsection{Γραφική απεικόνιση του καλύτερου regression μοντέλου}

Στο Σχήμα~\ref{fig:best_lr_fit} φαίνεται η προσαρμογή του καλύτερου μοντέλου παλινδρόμησης στο training set.

\begin{figure}[htbp]
    \centering
    \includegraphics[width=0.95\textwidth]{lab6_2.png}
    \caption{Προσαρμογή του καλύτερου regression μοντέλου στο training set.}
    \label{fig:best_lr_fit}
\end{figure}

Το μοντέλο δεν προσπαθεί να ακολουθήσει όλες τις αιχμές της χρονοσειράς, κάτι που θα οδηγούσε σε υπερπροσαρμογή. Αντίθετα, συλλαμβάνει τη βασική περιοδική δομή και τη μέση μεταβολή του σήματος. Αυτό είναι επιθυμητό, επειδή οι απότομες αιχμές αποτελούν κυρίως τυχαίες διακυμάνσεις ή έκτακτα γεγονότα και δεν είναι εύκολο να προβλεφθούν από ένα απλό regression μοντέλο.

Στο Σχήμα~\ref{fig:best_lr_forecast} φαίνεται η πρόβλεψη του $LR3$ στο test interval μαζί με τα $95\%$ prediction intervals.

\begin{figure}[htbp]
    \centering
    \includegraphics[width=0.95\textwidth]{lab6_3.png}
    \caption{Πρόβλεψη του καλύτερου regression μοντέλου με $95\%$ prediction intervals.}
    \label{fig:best_lr_forecast}
\end{figure}

Τα prediction intervals καλύπτουν σχεδόν όλες τις πραγματικές τιμές. Η τελική κάλυψη είναι $97.87\%$, ελαφρώς υψηλότερη από το ονομαστικό $95\%$. Αυτό δείχνει ότι τα διαστήματα είναι λίγο συντηρητικά, αλλά παραμένουν χρήσιμα. Η μεγάλη αιχμή κοντά στην ώρα $1239$ δεν προβλέπεται σημειακά και βρίσκεται εκτός του διαστήματος πρόβλεψης. Αυτό είναι φυσιολογικό: ένα τέτοιο spike είναι δύσκολο να προβλεφθεί από deterministic terms, αλλά μπορεί να καλυφθεί από την εκτιμημένη αβεβαιότητα, αν και όταν αυτό γίνει υπερβολικά μεγάλο, τότε είναι αναμενόμενο να ξεφύγει από τα προδιαγεγραμμένα διαστήματα πρόβλεψης.

\section{Αποτελέσματα φίλτρων}

\subsection{Σύγκριση υποψήφιων φίλτρων}

Τα αποτελέσματα των πέντε φίλτρων φαίνονται στον Πίνακα~\ref{tab:filter_results}.

\begin{table}[htbp]
\centering
\small
\caption{Αποτελέσματα υποψήφιων μοντέλων φίλτρων.}
\label{tab:filter_results}
\begin{tabular}{|P{5.2cm}|P{1.5cm}|P{2.1cm}|P{1.8cm}|P{1.8cm}|P{2.1cm}|}
\hline
\textbf{Μοντέλο} & \textbf{Τάξη} & \textbf{Residual Std} & \textbf{MAE} & \textbf{RMSE} & \textbf{Median CRPS} \\
\hline
F1: $\Delta_{24}$ & 24 & 904.672 & 475.487 & 734.967 & 310.531 \\
\hline
F2: $\Delta_1\Delta_{24}$ & 25 & 1110.336 & 753.587 & 1043.372 & 1338.159 \\
\hline
F3: $\Delta_{168}$ & 168 & 955.139 & 616.918 & 864.900 & 296.036 \\
\hline
F4: $\Delta_1\Delta_{168}$ & 169 & 1176.400 & 970.518 & 1219.322 & 1402.320 \\
\hline
F5: $\Delta_{24}\Delta_{168}$ & 192 & 1267.724 & 832.639 & 1308.248 & 456.351 \\
\hline
\end{tabular}
\end{table}

Το καλύτερο φίλτρο με βάση το median CRPS είναι:
\[
F3:\quad \Delta_{168}.
\]
Δηλαδή, το καλύτερο φίλτρο είναι η εβδομαδιαία διαφοροποίηση:
\[
\Delta_{168}Y_t = Y_t - Y_{t-168}.
\]

Αξίζει να σημειωθεί ότι το $F1=\Delta_{24}$ έχει καλύτερο $MAE$ και $RMSE$ από το $F3$, αλλά χειρότερο median CRPS. Επομένως, αν εξετάζαμε μόνο τις σημειακές προβλέψεις, το ημερήσιο φίλτρο θα ήταν ισχυρότερος υποψήφιος. Όμως, επειδή στην άσκηση ζητούνται και prediction intervals, το CRPS είναι πιο ολοκληρωμένο κριτήριο. Με βάση αυτό, το εβδομαδιαίο φίλτρο δίνει καλύτερη συνολική πιθανοτική πρόβλεψη.

Τα φίλτρα που περιλαμβάνουν επιπλέον $\Delta_1$ έχουν σαφώς χειρότερη επίδοση. Αυτό δείχνει ότι η πρώτη διαφοροποίηση πιθανότατα οδηγεί σε over-differencing. Με άλλα λόγια, αφαιρεί περισσότερη δομή από όση χρειάζεται, αυξάνει τη μεταβλητότητα των residuals και τελικά οδηγεί σε χειρότερες προβλέψεις και ευρύτερα ή λιγότερο αποδοτικά prediction intervals.

\subsection{Διαγνωστικά residuals του καλύτερου φίλτρου}

Στο Σχήμα~\ref{fig:best_filter_residuals} φαίνονται τα κεντραρισμένα residuals του φίλτρου $\Delta_{168}$.

\begin{figure}[htbp]
    \centering
    \includegraphics[width=0.95\textwidth]{lab6_4.png}
    \caption{Κεντραρισμένα residuals του καλύτερου φίλτρου $F3=\Delta_{168}$.}
    \label{fig:best_filter_residuals}
\end{figure}

Η φιλτραρισμένη σειρά κινείται γενικά γύρω από το μηδέν, κάτι που είναι επιθυμητό για ένα μοντέλο της μορφής:
\[
F(B)Y_t = c + \epsilon_t.
\]
Ωστόσο, παραμένουν έντονες αιχμές. Οι αιχμές αυτές δείχνουν ότι η χρονοσειρά δεν μπορεί να περιγραφεί πλήρως από ένα απλό εποχικό φίλτρο. Παρόλα αυτά, το φίλτρο $\Delta_{168}$ καταφέρνει να εκμεταλλευτεί την εβδομαδιαία επανάληψη του φορτίου και να δώσει την καλύτερη πιθανοτική επίδοση μεταξύ των φίλτρων.

\subsection{Πρόβλεψη με το καλύτερο φίλτρο}

Στο Σχήμα~\ref{fig:best_filter_forecast} φαίνεται η πρόβλεψη του καλύτερου φίλτρου στο test interval.

\begin{figure}[htbp]
    \centering
    \includegraphics[width=0.95\textwidth]{lab6_5.png}
    \caption{Πρόβλεψη του καλύτερου φίλτρου $F3=\Delta_{168}$ με $95\%$ prediction intervals.}
    \label{fig:best_filter_forecast}
\end{figure}

Το φίλτρο ακολουθεί πιο απότομα τις τιμές της προηγούμενης εβδομάδας και γι' αυτό παρουσιάζει πιο έντονη τοπική μεταβλητότητα από το regression. Τα prediction intervals είναι αρκετά πλατιά ώστε να καλύπτουν τις περισσότερες πραγματικές τιμές, με τελική κάλυψη $97.87\%$. Το εύρος αυτό είναι αναμενόμενο, επειδή στα φίλτρα η αβεβαιότητα διαδίδεται μέσω της αντίστροφης αναπαράστασης του φίλτρου.

\section{Τελική σύγκριση επιλεγμένων μοντέλων}

\subsection{Αριθμητική σύγκριση}

Στον Πίνακα~\ref{tab:selected_comparison} συγκρίνονται τα δύο τελικά επιλεγμένα μοντέλα, δηλαδή το καλύτερο regression και το καλύτερο filter.

\begin{table}[htbp]
\centering
\caption{Τελική σύγκριση επιλεγμένων μοντέλων.}
\label{tab:selected_comparison}
\begin{tabular}{|P{8cm}|P{2cm}|P{2cm}|P{2cm}|P{2.3cm}|}
\hline
\textbf{Επιλεγμένο μοντέλο} & \textbf{MAE} & \textbf{RMSE} & \textbf{Coverage} & \textbf{Median CRPS} \\
\hline
LR3: linear trend + daily 24 + weekly 168 & 419.878 & 675.712 & 97.87\% & 209.130 \\
\hline
F3: $\Delta_{168}$ & 616.918 & 864.900 & 97.87\% & 296.036 \\
\hline
\end{tabular}
\end{table}

Το regression μοντέλο $LR3$ υπερέχει σε όλες τις βασικές μετρικές σφάλματος και στο median CRPS, ενώ έχει την ίδια κάλυψη prediction interval με το καλύτερο φίλτρο. Συνεπώς, αποτελεί την καλύτερη συνολική επιλογή για τα συγκεκριμένα δεδομένα και τον συγκεκριμένο ορίζοντα πρόβλεψης.

Η ίση κάλυψη $97.87\%$ δείχνει ότι και τα δύο μοντέλα δίνουν επαρκώς συντηρητικά $95\%$ intervals. Όμως, το regression έχει χαμηλότερο CRPS, άρα πετυχαίνει καλύτερο συμβιβασμό μεταξύ ακρίβειας σημειακής πρόβλεψης και εύρους αβεβαιότητας.

\subsection{Κοινή γραφική σύγκριση με διαστήματα πρόβλεψης}

Στο Σχήμα~\ref{fig:common_forecast_intervals} παρουσιάζεται η κοινή γραφική σύγκριση των δύο τελικά επιλεγμένων μοντέλων. Στο ίδιο διάγραμμα φαίνονται οι πραγματικές τιμές του forecast interval, οι σημειακές προβλέψεις της γραμμικής παλινδρόμησης και του φίλτρου, καθώς και τα αντίστοιχα $95\%$ prediction intervals.

\begin{figure}[htbp]
    \centering
    \includegraphics[width=0.95\textwidth]{lab6_8.png}
    \caption{Κοινή σύγκριση των επιλεγμένων μοντέλων με $95\%$ prediction intervals.}
    \label{fig:common_forecast_intervals}
\end{figure}

Το διάγραμμα επιβεβαιώνει οπτικά τα αριθμητικά αποτελέσματα του Πίνακα~\ref{tab:selected_comparison}. Το regression μοντέλο δίνει πιο ομαλή πρόβλεψη και στενότερη περιοδική δομή, ενώ το filter μοντέλο εμφανίζει εντονότερες τοπικές διακυμάνσεις, επειδή βασίζεται στην εβδομαδιαία υστέρηση. Και τα δύο μοντέλα παρέχουν επαρκή κάλυψη των πραγματικών τιμών, όμως το regression μοντέλο επιτυγχάνει καλύτερο συνολικό συμβιβασμό μεταξύ ακρίβειας και αβεβαιότητας, όπως φαίνεται από το μικρότερο median CRPS.

\subsection{Σύγκριση σημειακών προβλέψεων}

Στο Σχήμα~\ref{fig:point_comparison} φαίνονται οι σημειακές προβλέψεις των δύο επιλεγμένων μοντέλων σε σχέση με τις πραγματικές τιμές.

\begin{figure}[htbp]
    \centering
    \includegraphics[width=0.95\textwidth]{lab6_6.png}
    \caption{Σύγκριση σημειακών προβλέψεων των επιλεγμένων μοντέλων.}
    \label{fig:point_comparison}
\end{figure}

Η παλινδρόμηση δίνει πιο ομαλή πρόβλεψη, επειδή βασίζεται σε deterministic αρμονικούς όρους. Αντίθετα, το φίλτρο δίνει πιο ευμετάβλητη πρόβλεψη, επειδή ουσιαστικά αξιοποιεί παρελθοντικές τιμές με βάση την εβδομαδιαία υστέρηση. Η ομαλή συμπεριφορά του regression αποδεικνύεται πιο κατάλληλη στο συγκεκριμένο forecast interval, επειδή μειώνει το μέσο σφάλμα χωρίς να χάνει τη βασική περιοδική δομή.

Στο Σχήμα~\ref{fig:absolute_errors} παρουσιάζεται η σύγκριση των απόλυτων σφαλμάτων.

\begin{figure}[h!]
    \centering
    \includegraphics[width=0.95\textwidth]{lab6_7.png}
    \caption{Σύγκριση απόλυτου σφάλματος των επιλεγμένων μοντέλων.}
    \label{fig:absolute_errors}
\end{figure}

Και τα δύο μοντέλα αποτυγχάνουν να προβλέψουν σημειακά τη μεγάλη αιχμή κοντά στην ώρα $1239$. Αυτό είναι αναμενόμενο, γιατί τέτοιες ακραίες τιμές δεν εξηγούνται εύκολα από απλή εποχικότητα ή γραμμική τάση. Ωστόσο, εκτός από αυτήν την αιχμή, το regression παρουσιάζει γενικά μικρότερα σφάλματα από το φίλτρο, κάτι που συμφωνεί με τα συνολικά $MAE$ και $RMSE$.

\subsection{Έξοδος του προγράμματος σε Octave}

Παρακάτω παρατίθεται η βασική έξοδος του προγράμματος από το Octave.

{\small
\begin{verbatim}
============================================================
Lab 6 - Basic Forecasting with Model Selection
============================================================
Loaded dataset: dataset.mat
Number of observations: 1247
Training observations: 1200
Forecast observations: 47
Forecast horizon: 1201 to 1247
Confidence level: 95%
============================================================
Linear Regression Candidate Models
============================================================
Model                                                        Params    Train MSE          MAE         RMSE     Med CRPS
LR1: constant + daily 24                                          3   521883.168      538.035      791.283      260.300
LR2: linear trend + daily 24                                      4   513540.892      458.657      710.656      232.094
LR3: linear trend + daily 24 + weekly 168                         6   498776.282      419.878      675.712      209.130
LR4: quadratic trend + daily 24 + weekly 168                      7   497531.338      437.483      702.799      215.583
LR5: linear trend + daily 24 + half-daily 12 + weekly 168         8   490046.575      402.893      644.796      235.488

Best Linear Regression model by median CRPS:
 LR3: linear trend + daily 24 + weekly 168
============================================================
Filter Candidate Models
============================================================
Model                               Order      Res Std          MAE         RMSE     Med CRPS
F1: Delta_24                           24      904.672      475.487      734.967      310.531
F2: Delta_1 Delta_24                   25     1110.336      753.587     1043.372     1338.159
F3: Delta_168                         168      955.139      616.918      864.900      296.036
F4: Delta_1 Delta_168                 169     1176.400      970.518     1219.322     1402.320
F5: Delta_24 Delta_168                192     1267.724      832.639     1308.248      456.351

Best Filter model by median CRPS:
 F3: Delta_168
============================================================
Final Selected Models Comparison
============================================================
Selected model                              MAE         RMSE     Coverage     Med CRPS
LR3: linear trend + daily 24 + weekly 168   419.878      675.712       97.87%      209.130
F3: Delta_168                              616.918      864.900       97.87%      296.036

Overall, the selected Linear Regression model has better median CRPS.
Saved figures as PNG files in the current directory.
Saved model comparison metrics to lab6_model_comparison_results.csv.
Execution completed successfully.
\end{verbatim}
}

\section{Συνολικό συμπέρασμα}

Στην παρούσα άσκηση εφαρμόστηκαν δύο βασικές κατηγορίες μεθόδων πρόβλεψης χρονοσειρών: γραμμική παλινδρόμηση και φίλτρα διαφοροποίησης. Για κάθε κατηγορία δοκιμάστηκαν πέντε θεωρητικά αιτιολογημένες παραλλαγές, ώστε η τελική επιλογή να μη βασιστεί σε ένα αυθαίρετο μοντέλο.

Στην κατηγορία της γραμμικής παλινδρόμησης, το καλύτερο μοντέλο ήταν το $LR3$, το οποίο περιλαμβάνει γραμμική τάση, ημερήσια περιοδικότητα $24$ ωρών και εβδομαδιαία περιοδικότητα $168$ ωρών. Η επιλογή αυτή είναι συνεπής με τη φύση των δεδομένων, αφού σε ωριαίες χρονοσειρές φόρτου είναι αναμενόμενο να εμφανίζονται ημερήσια και εβδομαδιαία πρότυπα.

Στην κατηγορία των φίλτρων, το καλύτερο μοντέλο ήταν το $F3=\Delta_{168}$, δηλαδή το φίλτρο εβδομαδιαίας διαφοροποίησης. Το αποτέλεσμα αυτό δείχνει ότι η εβδομαδιαία επανάληψη είναι σημαντική στη συγκεκριμένη χρονοσειρά. Τα φίλτρα που περιλάμβαναν επιπλέον πρώτη διαφοροποίηση είχαν χειρότερη συμπεριφορά, κάτι που υποδηλώνει ότι η επιπλέον διαφοροποίηση οδηγεί σε over-differencing.

Η τελική σύγκριση έδειξε ότι το regression μοντέλο $LR3$ υπερέχει του καλύτερου φίλτρου $F3$. Συγκεκριμένα, πέτυχε μικρότερο $MAE$, μικρότερο $RMSE$ και μικρότερο median CRPS, ενώ είχε την ίδια κάλυψη των $95\%$ prediction intervals. Επομένως, για τα συγκεκριμένα δεδομένα και τον συγκεκριμένο forecast horizon, η γραμμική παλινδρόμηση με ημερήσια και εβδομαδιαία εποχικότητα αποτελεί την καλύτερη από τις εξεταζόμενες επιλογές.

Τέλος, τα αποτελέσματα αναδεικνύουν και ένα γενικό συμπέρασμα για την πρόβλεψη φόρτου σε συστήματα αναμονής: οι ομαλές περιοδικές συνιστώσες μπορούν να προβλεφθούν ικανοποιητικά, όμως οι απότομες αιχμές παραμένουν δύσκολες. Για τον λόγο αυτό, η χρήση prediction intervals είναι απαραίτητη, επειδή εκφράζει την αβεβαιότητα που δεν μπορεί να αποτυπωθεί από τη σημειακή πρόβλεψη.

\newpage
\section{Παράρτημα Υλοποίησης σε Octave}

\subsection{Κώδικας Υλοποίησης}
\includepdf[pages=-]{lab6.pdf}

\subsection{Έξοδος Προγράμματος}
\includepdf[pages=-]{lab6_output.pdf}

\end{document}